
Multi-dimensional trustworthiness evaluation treats reliability, robustness, and fairness as independent properties assessed through specialized benchmarks. This study tests empirical orthogonality by computing cross-benchmark correlations across 20 large language models spanning three parameter size strata. Analysis reveals near-perfect positive correlations (r = 0.993--0.998, p < $2\times 10^{-17})$ across TrustfulQA (reliability), AdvBench (robustness), and BOLD (fairness) benchmarks, far exceeding hypothesized moderate effect sizes (r > 0.3) expected from shared root causes. Correlations persist within size strata (r > 0.98), indicating coupling is not attributable to model scale. Hierarchical clustering analysis yields perfectly stable groups (100\% bootstrap consistency over 1000 iterations) but poor separation (silhouette score = 0.274 < 0.5 threshold), establishing that the 3-benchmark design cannot validate failure mode taxonomy despite adequate statistical power for correlation detection. Results support two competing explanations: (1) benchmarks measure a unified construct rather than orthogonal dimensions, or (2) distinct dimensions exist but 3 benchmarks provide insufficient resolution. Current data cannot distinguish these hypotheses without experimental manipulation or expansion to 5--10 benchmarks. The tight coupling challenges independent evaluation paradigms and suggests multi-dimensional interventions may be more efficient than dimension-specific fixes, though intervention validation remains untested.
